%% kapa: exam=1 round=25 year=2014 period=2 form=A subject=accounting subject_ko=회계학 seq=5 number=45
%% source: Q-Net 감정평가사 기출문제 (cst00309) artlSeq=5101207 file=extracted/1008/5101207/2014년 제25회 감정평가사 1차시험 기출문제/2교시(A형).hwp
%% conversion: kordoc 4.24.51 -> tools/gen1.py (auto); review=proofread-sonnet; flags=none
\begin{question}{45}
유형자산에 관한 설명으로 옳은 것은?
\begin{choices}
  \item 유형자산은 다른 자산의 미래경제적효익을 얻기 위해 필요하더라도, 그 자체로의 직접적인 미래경제적효익을 얻을 수 없다면 인식할 수 없다.
  \item 유형자산이 경영진이 의도하는 방식으로 가동될 수 있으나 가동수준이 완전조업도 수준에 미치지 못하는 경우에 발생하는 원가는 유형자산 원가에 포함한다.
  \item 유형자산의 원가는 경영진이 의도하는 방식으로 자산을 가동하는 데 필요한 장소와 상태에 이르게 하는 데 직접 관련되는 원가를 포함하며, 해당 자산의 시험과정에서 생산된 시제품의 순매각금액은 차감한다.
  \item 건설이 시작되기 전에 건설용지를 주차장 용도로 사용함에 따라 획득한 수익은 유형자산의 원가에서 차감한다.
  \item 교환거래에 상업적 실질이 있는지 여부를 결정할 때 교환거래의 영향을 받는 영업 부분의 기업특유가치는 세전현금흐름을 반영하여야 한다.
\end{choices}
\end{question}
\answer{3}
